\section{Related Work}
\label{sec:relatedwork}

This section reviews the multi-agent and reinforcement-learning systems that established cooperative agent modelling as a practical tool for sports and business analytics (Subsection~\ref{subsec:magentsystems}), before examining why none of these systems can support a single, video-grounded, faithfulness-checked, multi-stakeholder cricket system such as the one proposed in this paper (Subsection~\ref{subsec:verificationgap}).

\subsection{Multi-Agent Modelling and Sports Analytics Systems}
\label{subsec:magentsystems}
Yeh \textit{et al.} \cite{ref1} established graph-based multi-agent trajectory modelling for team sports, combining a graph neural network with a variational recurrent network to obtain a permutation-equivariant model that outperformed fixed-ordering RNN and VRNN baselines on basketball and soccer tracking data. Teoh \cite{ref2} extended this direction with CausalTraj, a temporally causal, likelihood-based model that, unlike prior work optimized only for per-agent accuracy, explicitly targets joint scenario coherence, reporting the best recorded joint-metric (minJADE, minJFDE) performance on the NBA SportVU, Basketball-U and Football-U benchmarks while remaining competitive on standard per-agent metrics. Complementing these trajectory models, the FOOTPASS project \cite{ref3} released a multi-modal dataset pairing broadcast video with structured tactical context, and showed that grounding action-spotting models in this structured context roughly doubled event-detection accuracy compared with using video alone. At the environment level, Kurach \textit{et al.} \cite{ref4} built the Google Research Football Environment, an open-source, physics-based simulator supporting single- and multi-agent reinforcement learning, and reported baseline results for IMPALA, PPO and Ape-X DQN across benchmark scenarios of varying difficulty. Building on such environments, Brandao \textit{et al.} \cite{ref5} used reinforcement learning to let simulated soccer agents automatically develop distinct tactical roles, reaching a positive goal difference against fixed rule-based opponents purely through reward shaping rather than hand-coded strategy. None of these systems, however, release a mechanism for verifying a numeric output against a trusted evidence source, nor are they designed to serve more than one downstream audience from the same architecture.

\subsection{Agentic Systems and the Verification Gap}
\label{subsec:verificationgap}
A more recent line of work applies coordinating large-language-model agents, rather than purely numerical trajectory or reinforcement-learning agents, to real decision problems. Gupta \cite{ref6} proposed a four-agent supply chain framework, covering demand forecasting, logistics, inventory and risk monitoring, coordinated through a hybrid edge-cloud collaboration layer, and reported a 22\% reduction in decision latency and 16\% faster disruption recovery relative to centralized and rule-based baselines. Hazenberg \textit{et al.} \cite{ref7} benchmarked several multi-agent reinforcement learning strategies for dynamic pricing across a simulated supply chain, showing that agent behaviour, and the resulting market outcome, depends heavily on how directly competing agents are trained against one another rather than on the learning algorithm alone. Closest to the present work, Bhatnagar \cite{ref8} introduced FanCric, a multi-agent system built from coordinating LLM agents that individually research player form, historical performance and match conditions before jointly selecting a fantasy cricket team, and showed this coordinated approach outperformed both a single-prompt baseline and aggregated human fan selections. While these agentic frameworks confirm that coordinating several specialized agents produces more reliable real-world decisions than a single monolithic model, none of them begins from raw video evidence, none verifies its agents' numeric claims against a shared, sanctioned source of truth before a report reaches the user, and none is designed to route a single underlying analysis to more than one type of stakeholder. No existing study combines video-grounded shot classification, retrieval-backed statistical evidence, dynamically routed multi-stakeholder agents and a dedicated faithfulness-checking gate within a single cricket analytics system, a gap this paper addresses through CricketMind.